\section{Exercise 3}
\textit{Let $F$ be a vector field on $\mathbb{R}^3$ given by:}
\[F(x,y,z)  = (yz, 2xz, e^{xy})\]

\begin{itemize}
    \item Does a potential function $\phi$ exist on $\mathbb{R}^3$, such that $\nabla\phi = F$
    \item Use Stokes theorem to evaluate $S=\oint_C F \cdot dr$, where C is the circle (oriented clockwise, viewed from above): $C = \{(x,y,z) \in \mathbf{R}^3 | x^2 + y^2 = 1, z=5\}$
\end{itemize}

To see if a potential function exists to yield the vector itself. The
potential is a scalar function. If the local rotation \textit{curl} $=0$ then
a potential function $\phi$ exists:
Given:
\[\begin{cases}
    P(x,y,z) = yz\\
    Q(x,y,z) = 2xz\\
    R(x,y,z) = e^{xy}
\end{cases}\]
\[\nabla \times F = 
\begin{pmatrix}
\frac{\partial R}{\partial y} - \frac{\partial Q}{\partial z}\\
\frac{\partial P}{\partial z} - \frac{\partial R}{\partial x}\\
\frac{\partial Q}{\partial x} - \frac{\partial P}{\partial y}
\end{pmatrix} 
=
\begin{pmatrix}
e^{xy}-2x\\
y-ye^{xy}\\
2z-z
\end{pmatrix} 
=
\begin{pmatrix}
x(e^{xy}-2)\\
y(1-e^{xy})\\
z
\end{pmatrix}
\]
Since there is local rotation, as in the curl is $\neq0$ , no function $\phi$
exists on $\mathbb{R}^2$.\\
Next it is possible to evaluate the integral $S$ using stokes theorem,
which states that a line closed integral on this form can be expressed as a 
double integal of the curl. The addition of the factor $n$, the normal
vector to the surface of the circle (or disk), which is perpendicular to its 
surface. And because we're oriented closewise, and viewed downwards (from above), that
normal vector becomes $n=(0,0,-1)$
\[S=\oint_C F \cdot dr = \iint_s (\nabla \times \mathbf{F}) \  n \ dS\]
Evaluation of the integrant:
\[(\nabla \times \mathbf{F}) n = 
\begin{pmatrix}
x(e^{xy}-2)\\
y(1-e^{xy})\\
z
\end{pmatrix}
\cdot 
\begin{pmatrix}
    0\\
    0\\
    -1
\end{pmatrix}
= x(e^{xy}-2)(0) + y(1-e^{xy})(0) + z(-1) 
=-5 \]
In this case, the integration is around of a circle, this makes the bounds: 
\[r^2 = 1,\qquad 0\le r \le1, \qquad 0\le \theta \le 2\pi\]
So outer integral sums every angle $\theta$ for every $r$, leading to:
\[\int_{0}^{1}\int_{0}^{2\pi}-5r \ d\theta \ dr \Rightarrow -5r \int_{0}^{1}\int_{0}^{2\pi} \ d\theta \ dr\] 
Inner integral:
\[-5r\int_{0}^{2\pi} \ d\theta = \left[-5r \ \theta\right]_0^{2\pi} = -10\pi r\]
Outer integral:
\[-10\pi \ \int_{0}^{1} r \ dr \Rightarrow \left[-10\pi \ \frac{1}{2}r^2\right]_0^1 = \left[-5\pi \ r^2 \right]_0^1 = -5\pi \]
\[\]
